% Séquence 6 : Proportionnalité (Partie 1)
% Durée estimée : 4 séances
\setseqtitle{Proportionnalité (Partie 1)}
\chapter{Proportionnalité — Partie 1}
\label{chap:seq06-proportionnalite-1}

% ============================================================
% SÉANCE 1 : Découverte de la proportionnalité
% ============================================================
\seance

\begin{objectifsbox}
	\textbf{Notions :} situation de proportionnalité, tableau de proportionnalité, coefficient de proportionnalité\\
	\textbf{Compétences :} reconnaître une situation de proportionnalité, compléter un tableau de proportionnalité\\
	\textbf{Applications :} prix, recettes, situations de la vie courante\\
	\textbf{Contrôle (prévisionnel) :} vendredi 9 janvier 2026
\end{objectifsbox}

\section{Introduction : qu'est-ce que la proportionnalité ?}

Dans la vie quotidienne, on rencontre souvent des situations où deux grandeurs sont liées de manière particulière : quand l'une double, l'autre double aussi ; quand l'une triple, l'autre triple aussi.

\vspace{0.5cm}
\begin{activitybox}[Situation 1 : Au marché]
	Chez le marchand de fruits, les pommes coûtent 2~€ le kilogramme.

	\begin{enumerate}
		\item Complète le tableau suivant :
		\begin{center}
			\renewcommand{\arraystretch}{1.5}
			\begin{tabular}{|c|c|c|c|c|c|}
				\hline
				\textbf{Masse de pommes (kg)} & 1 & 2 & 3 & 5 & 10 \\
				\hline
				\textbf{Prix (€)} & \solution{2} & \solution{4} & \solution{6} & \solution{10} & \solution{20} \\
				\hline
			\end{tabular}
		\end{center}

		\solution{
			\textbf{Solution :} Les pommes coûtent 2~€ le kg, donc :
			\begin{itemize}
				\item 1 kg coûte $1 \times 2 = 2$~€
				\item 2 kg coûtent $2 \times 2 = 4$~€
				\item 3 kg coûtent $3 \times 2 = 6$~€
				\item 5 kg coûtent $5 \times 2 = 10$~€
				\item 10 kg coûtent $10 \times 2 = 20$~€
			\end{itemize}
		}

		\item Que remarques-tu entre la masse et le prix ?

		\solution[8cm]{Le prix est toujours égal à 2 fois la masse. Pour trouver le prix, on multiplie toujours la masse par 2.}

		\item Si on achète 2 fois plus de pommes, que se passe-t-il pour le prix ?

		\solution[6cm]{Le prix est aussi 2 fois plus grand. Par exemple, si on passe de 5 kg (10~€) à 10 kg, le prix passe de 10~€ à 20~€.}
	\end{enumerate}
\end{activitybox}

\begin{activitybox}[Situation 2 : L'âge et la taille]
	Voici un tableau donnant l'âge et la taille de Paul :

	\begin{center}
		\renewcommand{\arraystretch}{1.5}
		\begin{tabular}{|c|c|c|c|c|}
			\hline
			\textbf{Âge (années)} & 2 & 4 & 8 & 12 \\
			\hline
			\textbf{Taille (cm)} & 85 & 102 & 128 & 152 \\
			\hline
		\end{tabular}
	\end{center}

	\begin{enumerate}
		\item Quand l'âge double (de 2 à 4 ans), la taille double-t-elle ?

		\solution[4cm]{
			Non. Vérifions : quand Paul a 2 ans, il mesure 85 cm. Si la taille doublait, à 4 ans il devrait mesurer $85 \times 2 = 170$ cm. Or, le tableau indique qu'à 4 ans, il mesure 102 cm. Donc $170 \neq 102$, la taille ne double pas.
		}

		\item Cette situation est-elle similaire à celle des pommes ?

		\solution[4cm]{
			Non, ce n'est pas similaire. Dans la situation des pommes, quand la masse double, le prix double aussi. Ici, quand l'âge double, la taille ne double pas. Il n'y a pas de nombre fixe permettant de passer de l'âge à la taille : ce n'est pas proportionnel.
		}
	\end{enumerate}
\end{activitybox}

\section{Définition de la proportionnalité}

\begin{definitionbox}[Situation de proportionnalité]
	Deux grandeurs sont \textbf{proportionnelles} si on peut passer de l'une à l'autre en multipliant toujours par le \solution[3cm]{même nombre non nul}.

	Ce nombre s'appelle le \textbf{coefficient de proportionnalité}.
\end{definitionbox}

\begin{examplebox}
	\textbf{Situation de proportionnalité :} Le prix des pommes est proportionnel à la masse.
	\begin{itemize}
		\item Pour passer de la masse au prix, on multiplie toujours par \solution{2}.
		\item Le coefficient de proportionnalité est \solution{2}.
	\end{itemize}

	\textbf{Situation de non proportionnalité :} L'âge et la taille de Paul ne sont pas proportionnels.
	\begin{itemize}
		\item Il n'existe pas de nombre fixe permettant de passer de l'âge à la taille.
	\end{itemize}
\end{examplebox}

\begin{remarkbox}[Attention aux pièges !]
	Il existe des situations qui \textbf{ressemblent} à de la proportionnalité mais qui n'en sont pas :

	\textbf{1. Les situations avec une valeur fixe (\og point de départ\fg)}:\\
	Si on paye un abonnement ou une prise en charge fixe, ce n'est pas proportionnel.
	\begin{itemize}
		\item \textit{Exemple :} Un taxi demande 5~€ de prise en charge puis 2~€ par km.
		\item Pour 10 km : $\solution[2cm]{5 + (2 \times 10)} = \solution{25}$~€.
		\item Pour 20 km : $\solution[2cm]{5 + (2 \times 20)} = \solution{45}$~€.
		\item Quand la distance double, \solution[7cm]{le prix ne double pas} (25~€ $\times 2 = 50 \neq 45$).
	\end{itemize}

	\textbf{2. L'intuition de l'addition (\og+2 donc +2\fg)}:\\
	Ce n'est pas parce qu'on ajoute 2 à la première grandeur qu'on ajoute 2 à la deuxième !
	\begin{itemize}
		\item \textit{Exemple :} Un stylo coûte 2~€.
		\item Si j'achète 3 stylos (c'est $1 + 2$ stylos), je ne paye pas $2 + 2 = 4$~€.
		\item Je paye $3 \times 2 = \solution{6}$~€.
		\item On ne peut pas juste additionner \og 2 \fg au prix parce qu'on a ajouté \og 2 \fg{} à la quantité.
	\end{itemize}
\end{remarkbox}

\clearpage
% ============================================================
% SÉANCE 2 : Reconnaître une situation de proportionnalité
% ============================================================
\seance

\section{Comment reconnaître une situation de proportionnalité ?}

\begin{methodebox}[Méthode 1 : Vérifier les quotients]
	Pour vérifier si un tableau représente une situation de proportionnalité :
	\begin{enumerate}
		\item Calculer le quotient $\dfrac{\text{2ème ligne}}{\text{1ère ligne}}$ pour chaque colonne
		\item Si tous les quotients sont \solution[2cm]{égaux}, c'est proportionnel
		\item Ce quotient commun est le \solution[4cm]{coefficient de proportionnalité}
	\end{enumerate}
\end{methodebox}

\begin{examplebox}[Application de la méthode]
	Ce tableau est-il un tableau de proportionnalité ?

	\begin{center}
		\renewcommand{\arraystretch}{1.5}
		\begin{tabular}{|c|c|c|c|c|}
			\hline
			\textbf{Nombre de cahiers} & 2 & 5 & 8 & 10 \\
			\hline
			\textbf{Prix (€)} & 3 & 7,5 & 12 & 15 \\
			\hline
		\end{tabular}
	\end{center}

	\textbf{Calcul des quotients :}
	\begin{itemize}
		\item $\dfrac{3}{2} = \solution{1,5}$
		\item $\dfrac{7,5}{5} = \solution{1,5}$
		\item $\dfrac{12}{8} = \solution{1,5}$
		\item $\dfrac{15}{10} = \solution{1,5}$
	\end{itemize}

	\textbf{Conclusion :} Tous les quotients sont égaux à 1,5, donc c'est un tableau de \solution[4cm]{proportionnalité}.

	Le coefficient de proportionnalité est \solution{1,5}. Cela signifie qu'un cahier coûte \solution{1,50}~€.
\end{examplebox}

\begin{methodebox}[Méthode 2 : Vérifier les propriétés de linéarité]
	Dans un tableau de proportionnalité :
	\begin{itemize}
		\item Si on \textbf{multiplie} une valeur de la 1ère ligne par un nombre, la valeur correspondante de la 2ème ligne est aussi \solution[3cm]{multipliée} par ce nombre.
		\item Si on \textbf{additionne} deux colonnes, on obtient une nouvelle colonne \solution[2cm]{valide}.
	\end{itemize}
\end{methodebox}

\begin{exercisebox}[Exercices d'entraînement]
	\textbf{Exercice 1 :} Ces tableaux sont-ils des tableaux de proportionnalité ? Justifie.

	\textbf{Tableau A :}
	\begin{center}
		\begin{tabular}{|c|c|c|c|}
			\hline
			4 & 8 & 10 & 14 \\
			\hline
			10 & 20 & 25 & 35 \\
			\hline
		\end{tabular}
	\end{center}

	\solution[10cm]{
		\textbf{Réponse : Oui}, c'est un tableau de proportionnalité.

		\textbf{Justification :} Calculons les quotients :
		\begin{itemize}
			\item $\dfrac{10}{4} = 2,5$
			\item $\dfrac{20}{8} = 2,5$
			\item $\dfrac{25}{10} = 2,5$
			\item $\dfrac{35}{14} = 2,5$
		\end{itemize}
		Tous les quotients sont égaux à 2,5. C'est donc un tableau de proportionnalité de coefficient 2,5.
	}

	\textbf{Tableau B :}
	\begin{center}
		\begin{tabular}{|c|c|c|c|}
			\hline
			2 & 4 & 6 & 8 \\
			\hline
			5 & 9 & 13 & 17 \\
			\hline
		\end{tabular}
	\end{center}

	\solution[10cm]{
		\textbf{Réponse : Non}, ce n'est pas un tableau de proportionnalité.

		\textbf{Justification :} Calculons les quotients :
		\begin{itemize}
			\item $\dfrac{5}{2} = 2,5$
			\item $\dfrac{9}{4} = 2,25$
		\end{itemize}
		Les quotients ne sont pas égaux ($2,5 \neq 2,25$). Ce n'est donc pas un tableau de proportionnalité.
	}
\end{exercisebox}

\clearpage
% ============================================================
% SÉANCE 3 : Le coefficient de proportionnalité
% ============================================================
\seance

\section{Le coefficient de proportionnalité}

\begin{definitionbox}[Coefficient de proportionnalité]
	Dans une situation de proportionnalité, le \textbf{coefficient de proportionnalité} est le nombre par lequel on \solution[3cm]{multiplie} les valeurs de la première grandeur pour obtenir les valeurs de la deuxième grandeur.

	\textbf{Formule :} Si $a$ et $b$ sont deux valeurs correspondantes :
	$$\text{Coefficient} = \dfrac{b}{a}$$
\end{definitionbox}

\begin{remarkbox}[Sens du coefficient]
	Le coefficient de proportionnalité dépend du sens dans lequel on lit le tableau :
	\begin{itemize}
		\item De la 1ère ligne vers la 2ème ligne : coefficient $k$
		\item De la 2ème ligne vers la 1ère ligne : coefficient $\dfrac{1}{k}$ (l'inverse)
	\end{itemize}
\end{remarkbox}

\begin{examplebox}
	\begin{center}
		\renewcommand{\arraystretch}{1.5}
		\begin{tabular}{|c|c|c|c|}
			\hline
			\textbf{Distance (km)} & 50 & 100 & 150 \\
			\hline
			\textbf{Temps (h)} & 1 & 2 & 3 \\
			\hline
		\end{tabular}
	\end{center}

	\begin{itemize}
		\item Pour passer de la distance au temps : coefficient = $\dfrac{1}{50} = \solution{0,02}$
		\item Pour passer du temps à la distance : coefficient = $\dfrac{50}{1} = \solution{50}$
	\end{itemize}

	\textbf{Interprétation :} La vitesse est de \solution{50} km/h.
\end{examplebox}

\section{Compléter un tableau de proportionnalité}

\begin{methodebox}[Compléter un tableau]
	Pour compléter un tableau de proportionnalité :
	\begin{enumerate}
		\item Trouver le coefficient de proportionnalité à partir d'une colonne complète
		\item Multiplier chaque valeur connue par ce coefficient pour trouver la valeur manquante
	\end{enumerate}

	\textbf{Ou utiliser les propriétés de linéarité :}
	\begin{itemize}
		\item Multiplier une colonne par un nombre
		\item Additionner deux colonnes
	\end{itemize}
\end{methodebox}

\begin{examplebox}[Compléter avec le coefficient]
	Complète ce tableau de proportionnalité :

	\begin{center}
		\renewcommand{\arraystretch}{1.5}
		\begin{tabular}{|c|c|c|c|c|}
			\hline
			\textbf{Nombre de croissants} & 4 & 7 & 10 & 15 \\
			\hline
			\textbf{Prix (€)} & 3,60 & \solution{6,30} & \solution{9} & \solution{13,50} \\
			\hline
		\end{tabular}
	\end{center}

	\textbf{Solution :}
	\begin{enumerate}
		\item Coefficient = $\dfrac{3,60}{4} = \solution{0,90}$~€ par croissant
		\item $7 \times 0,90 = \solution{6,30}$~€
		\item $10 \times 0,90 = \solution{9}$~€
		\item $15 \times 0,90 = \solution{13,50}$~€
	\end{enumerate}
\end{examplebox}

\clearpage
% ============================================================
% SÉANCE 4 : Propriétés de linéarité
% ============================================================
\seance

\section{Propriétés de linéarité}

\begin{proprietebox}[Propriétés fondamentales]
	Dans un tableau de proportionnalité :

	\textbf{Propriété 1 (multiplicative) :} Si on multiplie une colonne par un nombre, on obtient une nouvelle colonne du tableau.

	\textbf{Propriété 2 (additive) :} Si on additionne deux colonnes, on obtient une nouvelle colonne du tableau.
\end{proprietebox}

\begin{examplebox}[Utilisation de la propriété multiplicative]
	\begin{center}
		\renewcommand{\arraystretch}{1.5}
		\begin{tabular}{|c|c|c|c|}
			\hline
			\textbf{Quantité de farine (g)} & 200 & ? & 500 \\
			\hline
			\textbf{Quantité de sucre (g)} & 80 & 120 & ? \\
			\hline
		\end{tabular}
	\end{center}

	\textbf{Pour trouver la quantité de farine correspondant à 120 g de sucre :}
	\begin{itemize}
		\item On sait que 200 g de farine $\rightarrow$ 80 g de sucre
		\item $120 = 80 \times 1,5$, donc on multiplie aussi 200 par \solution{1,5}
		\item $200 \times 1,5 = \solution{300}$ g de farine
	\end{itemize}

	\textbf{Pour trouver la quantité de sucre correspondant à 500 g de farine :}
	\begin{itemize}
		\item Coefficient = $\dfrac{80}{200} = \solution{0,4}$
		\item $500 \times 0,4 = \solution{200}$ g de sucre
	\end{itemize}
\end{examplebox}

\begin{examplebox}[Utilisation de la propriété additive]
	\begin{center}
		\renewcommand{\arraystretch}{1.5}
		\begin{tabular}{|c|c|c|c|}
			\hline
			\textbf{Nombre de personnes} & 4 & 6 & 10 \\
			\hline
			\textbf{Quantité de riz (g)} & 300 & 450 & ? \\
			\hline
		\end{tabular}
	\end{center}

	\textbf{Méthode :} $10 = 4 + 6$, donc :
	$$\text{Quantité pour 10 personnes} = 300 + 450 = \solution{750} \text{ g}$$
\end{examplebox}

\section{Applications concrètes}

\subsection{Les recettes de cuisine}

\begin{applicationbox}[Adapter une recette]
	Une recette de gâteau pour 4 personnes nécessite :
	\begin{itemize}
		\item 200 g de farine
		\item 120 g de sucre
		\item 2 œufs
		\item 100 g de beurre
	\end{itemize}

	\textbf{Question :} Quelles quantités faut-il pour 6 personnes ?

	\textbf{Solution :}

	\textit{Coefficient : $\dfrac{6}{4} = 1,5$. Pour 6 personnes, il faut multiplier toutes les quantités par 1,5.}

	\begin{center}
		\renewcommand{\arraystretch}{1.5}
		\begin{tabular}{|c|c|c|}
			\hline
			\textbf{Ingrédient} & \textbf{Pour 4 personnes} & \textbf{Pour 6 personnes} \\
			\hline
			Farine (g) & 200 & \solution{$200 \times 1,5 = 300$} \\
			\hline
			Sucre (g) & 120 & \solution{$120 \times 1,5 = 180$} \\
			\hline
			Œufs & 2 & \solution{$2 \times 1,5 = 3$} \\
			\hline
			Beurre (g) & 100 & \solution{$100 \times 1,5 = 150$} \\
			\hline
		\end{tabular}
	\end{center}
\end{applicationbox}

\begin{remarkbox}[Pour s'entra\^{i}ner]
	\textbf{Cahier Transmath 6e - Chapitre 16 (La proportionnalit\'e) :}
	\begin{itemize}
		\item Fiche 114 (p.125) \og Proportionnalit\'e et lin\'earit\'e\fg{} : ex. 20 \`a 25
		\item Fiche 115 (p.126) \og Retour \`a l'unit\'e\fg{} : ex. 26 \`a 32
		\item Fiche 116 (p.127) \og Repr\'esentations de la proportionnalit\'e\fg{} : ex. 33 \`a 41
	\end{itemize}
\end{remarkbox}

\begin{remarkbox}[Pour aller plus loin]
	La proportionnalité a de nombreuses applications concrètes que nous étudierons dans la Partie 2 :
	\begin{itemize}
		\item Les \textbf{pourcentages} : soldes, réductions, statistiques
		\item Les \textbf{échelles} : cartes géographiques, plans de bâtiments, maquettes
	\end{itemize}
	Ces applications utilisent directement les notions de coefficient de proportionnalité et de tableaux de proportionnalité vues dans cette séquence.
\end{remarkbox}

\section{Bilan de la séquence}

\begin{bilanbox}
	\textbf{Ce que je dois savoir :}
	\begin{itemize}
		\item[$\checkmark$] Définir ce qu'est une situation de proportionnalité
		\item[$\checkmark$] Connaître la définition du coefficient de proportionnalité
		\item[$\checkmark$] Connaître les propriétés de linéarité (additive et multiplicative)
	\end{itemize}

	\textbf{Ce que je dois savoir faire :}
	\begin{itemize}
		\item[$\checkmark$] Reconnaître si un tableau est un tableau de proportionnalité
		\item[$\checkmark$] Calculer un coefficient de proportionnalité
		\item[$\checkmark$] Compléter un tableau de proportionnalité
		\item[$\checkmark$] Utiliser la proportionnalité dans des situations concrètes
	\end{itemize}
\end{bilanbox}

\begin{summarybox}[Formules et méthodes à retenir]
	\begin{center}
		\fbox{\parbox{14cm}{
			\textbf{PROPORTIONNALITÉ}

			\vspace{0.3cm}
			\textbf{Définition :} Deux grandeurs sont proportionnelles si on multiplie toujours par le même nombre pour passer de l'une à l'autre.

			\vspace{0.3cm}
			\textbf{Coefficient de proportionnalité :}
			$$k = \dfrac{\text{valeur de la 2ème grandeur}}{\text{valeur de la 1ère grandeur}}$$

			\vspace{0.3cm}
			\textbf{Vérification :} Calculer tous les quotients $\dfrac{\text{2ème ligne}}{\text{1ère ligne}}$. S'ils sont tous égaux, c'est proportionnel.

			\vspace{0.3cm}
			\textbf{Propriétés de linéarité :}
			\begin{itemize}
				\item On peut multiplier une colonne par un nombre
				\item On peut additionner deux colonnes
			\end{itemize}
		}}
	\end{center}
\end{summarybox}
